% !TeX root = ../../Book5.tex
% 纲要标题：### 2.4 特征标表
% 纲要位置：工作记录/计划纲要.md，第 4100 行。
% 纲要细目（写作范围）：
% * 循环群、二面体群和四元数群
% * S_3、S_4 与 A_4 的计算
% * 由表能读出什么、不能读出什么
% ---


\section{特征标表}
\label{b5:sec:0204}

计算特征标表时，先确定共轭类，再从具体表示获得候选行，用内积确认不可约性，最后用次数平方和检查是否遗漏。以下均取复表示，并在表中列出共轭类大小；这些大小正是计算特征标内积时各列的权重。


\subsection{循环群、二面体群与四元数群}
\label{b5:subsec:0204:01}

对 \(C_m=\langle r\rangle\)，群交换，所有不可约表示一维。令 \(\zeta=\exp(2\pi\mathrm i/m)\)，则
\[
	\chi_j(r^a)=\zeta^{ja},\qquad 0\le j,a<m.
\]
这些行互异，已经有 \(m\) 个，次数平方和为 \(m\)，所以列表完整。正交关系也可由等比求和 \(\sum_{a=0}^{m-1}\zeta^{(i-j)a}=m\delta_{ij}\) 直接看到。

\ProseParagraphBreak
令 \(D_{2n}=\langle r,s\mid r^n=s^2=1,\ srs=r^{-1}\rangle\)，其中 \(n\ge3\)，下标为群阶。它的二维表示可从旋转子群的特征空间寻找：\(r^n=1\) 使 \(r\) 的复作用矩阵可对角化。若 \(rv=\lambda v\)，则由 \(rs=sr^{-1}\) 得
\[
	r(sv)=sr^{-1}v=\lambda^{-1}sv.
\]
因此 \(s\) 把 \(\lambda\) 特征空间送到 \(\lambda^{-1}\) 特征空间。若两特征值不同，\(v,sv\) 线性无关，张成一个在 \(r,s\) 下都不变的二维空间；在这组基下，\(r\) 为对角矩阵，\(s\) 交换两个基向量。再由 \(\lambda^n=1\)，取 \(\zeta=\exp(2\pi\mathrm i/n)\)，就得到下面的候选表示：对 \(j\in\mathbb Z/n\mathbb Z\)，定义
\[
	\rho_j(r)=\begin{pmatrix}\zeta^j&0\\0&\zeta^{-j}\end{pmatrix},
	\qquad
	\rho_j(s)=\begin{pmatrix}0&1\\1&0\end{pmatrix}.
\]
矩阵满足群关系。若 \(2j\not\equiv0\pmod n\)，两个旋转特征值不同，而反射交换对应特征直线，所以表示不可约。其特征标为
\[
	\chi_j(r^a)=\zeta^{ja}+\zeta^{-ja} =2\cos(2\pi ja/n),\qquad \chi_j(r^as)=0.
\]
交换两个基向量给出 \(\rho_j\simeq\rho_{-j}\)。反过来，若 \(\rho_j\simeq\rho_\ell\)，则旋转矩阵的特征值集合相同，故 \(\ell\equiv\pm j\pmod n\)；除此之外没有重复。若 \(2j\equiv0\pmod n\)，旋转矩阵为标量矩阵，两条直线 \(\mathbb C(1,1)\)、\(\mathbb C(1,-1)\) 均不变，故 \(\rho_j\) 可约。奇数 \(n\) 时只有 \(j=0\) 属于这种情形，偶数 \(n\) 时则有 \(j=0,n/2\)。一维表示由 \(r\mapsto a,s\mapsto b\) 给出。关系 \(r^n=s^2=1\) 要求 \(a^n=b^2=1\)，而 \(srs=r^{-1}\) 要求 \(bab=a^{-1}\)。标量交换且 \(b^2=1\)，故后一式成为 \(a=a^{-1}\)，即 \(a^2=1\)。因此 \(b=\pm1\)，而 \(a\) 在奇数 \(n\) 时只能为 \(1\)，在偶数 \(n\) 时可为 \(\pm1\)。当 \(n\) 奇数时有两个一维表示和 \((n-1)/2\) 个二维不可约表示；当 \(n\) 偶数时有四个一维表示和 \((n-2)/2\) 个二维不可约表示。两种情况下次数平方和均为 \(2n\)，所以穷尽不可约类型。

\ProseParagraphBreak
旋转彼此交换，而 \(sr^as^{-1}=r^{-a}\)，所以旋转的共轭类为 \(\{r^a,r^{-a}\}\)，相同元素时只算一个。对反射，有
\[
	r^b(r^as)r^{-b}=r^{a+2b}s,\qquad s(r^as)s^{-1}=r^{-a}s.
\]
这两种共轭只使指数增加一个偶数。若 \(n\) 为奇数，\(2\) 在 \(\mathbb Z/n\mathbb Z\) 中可逆，\(a+2b\) 遍历全部指数，所以反射成一个类；若 \(n\) 为偶数，这些变换保持指数的奇偶性，而 \(a+2b\) 遍历同一奇偶性的全部指数，所以恰分成两类。这给出将上述公式排成完整特征标表所需的全部列。

\ProseParagraphBreak
四元数群 \(Q_8=\{\pm1,\pm\mathbf i,\pm\mathbf j,\pm\mathbf k\}\) 的四个一维表示来自商 \(Q_8/\{\pm1\}\simeq C_2\times C_2\)。寻找非一维不可约表示时，中心元素 \(-1\) 由 Schur 引理按标量作用，其平方为一，所以只能作用为 \(I\) 或 \(-I\)。若为 \(I\)，表示会经过上述交换商群，因而不可约时只能一维；所以所求表示中必有 \(\rho(-1)=-I\)。于是 \(\rho(\mathbf i)^2=-I\)，其特征值为 \(\pm\mathrm i\)。关系 \(\mathbf j\mathbf i\mathbf j^{-1}=\mathbf i^{-1}\) 使 \(\mathbf j\) 交换这两个特征空间。取一个 \(\mathrm i\) 特征向量 \(v\)，则 \(v,-\rho(\mathbf j)v\) 张成不变二维空间；利用 \(\rho(\mathbf j)^2=-I\)，在这组基下自然得到
\[
	\rho(\mathbf i)=\begin{pmatrix}\mathrm i&0\\0&-\mathrm i\end{pmatrix},
	\qquad
	\rho(\mathbf j)=\begin{pmatrix}0&1\\-1&0\end{pmatrix}
\]
这两个矩阵平方均为 \(-I\)，且反交换，满足四元数关系，确实定义一个表示；前一个的特征直线被后一个交换，故不可约。这里与二面体群的区别在于，交换特征空间的矩阵平方为 \(-I\)，而反射的平方为 \(I\)。次数平方和 \(4+4=8\)，列表完整。\Cref{b5:tab:d8-q8}将四元数群与正方形对称群的特征标并排比较。表中的一维特征标按生成元取值编号：对 \(a,b\in\{1,-1\}\)，在 \(D_8\) 中令 \(\lambda_{a,b}(r)=a\)、\(\lambda_{a,b}(s)=b\)；在 \(Q_8\) 中则选生成元 \(\mathbf i,\mathbf j\)，规定 \(\mathbf k=\mathbf i\mathbf j\)，令 \(\lambda_{a,b}(\mathbf i)=a\)、\(\lambda_{a,b}(\mathbf j)=b\)。两群中这些表示都把第二列的中心元素送到 \(1\)，按各自表列的代表顺序，取值均为 \((1,1,a,b,ab)\)。

\begin{table}[htbp]
	\centering
	\caption{八阶二面体群与四元数群的共同特征标数值}
	\label{b5:tab:d8-q8}
	\begin{tabular}{lrrrrr}
		\toprule
		\(D_8\) 代表 & \(1\)&\(r^2\)&\(r\)&\(s\)&\(rs\)\\
		\(Q_8\) 代表 & \(1\)&\(-1\)&\(\mathbf i\)&\(\mathbf j\)&\(\mathbf k\)\\
		类大小 & \(1\)&\(1\)&\(2\)&\(2\)&\(2\)\\
		\midrule
		\(\lambda_{1,1}\)&1&1&1&1&1\\
		\(\lambda_{1,-1}\)&1&1&1&-1&-1\\
		\(\lambda_{-1,1}\)&1&1&-1&1&-1\\
		\(\lambda_{-1,-1}\)&1&1&-1&-1&1\\
		\(\chi\)&2&-2&0&0&0\\
		\bottomrule
	\end{tabular}
\end{table}

\subsection{\texorpdfstring{\(S_3\)、\(S_4\) 与 \(A_4\) 的特征标表}{S3、S4 与 A4 的特征标表}}
\label{b5:subsec:0204:02}

对 \(S_3\)，共轭类由轮换型给出，大小为 \(1,3,2\)。平凡与符号表示给出两行；三点置换表示减去平凡直和项，得到次数二的表示，特征标为 \(2,0,-1\)。其范数平方为 \((4+2)/6=1\)，故不可约，且平方和 \(1+1+4=6\) 已穷尽，得到\cref{b5:tab:s3}。

\begin{table}[htbp]
	\centering
	\caption{\(S_3\) 的特征标表}\label{b5:tab:s3}
	\begin{tabular}{lrrr}
		\toprule
		代表&\(1\)&\((12)\)&\((123)\)\\
		类大小&1&3&2\\
		\midrule
		\(1\)&1&1&1\\
		\(\mathrm{sgn}\)&1&-1&1\\
		\(\chi\)&2&0&-1\\
		\bottomrule
	\end{tabular}
\end{table}

由 \(S_3\) 的表可以实际核对列正交。对换位这一列，遍历全部不可约行，得到
\[
	1^2+(-1)^2+0^2=2=|C_{S_3}((12))|.
\]
换位与三轮换两列的内积为 \(1\cdot1+(-1)\cdot1+0\cdot(-1)=0\)。这里按不可约类型求和，未乘类大小；行正交则对共轭类求和，并以类大小加权。

\ProseParagraphBreak
对 \(S_4\)，五种轮换型的类大小为 \(1,6,3,8,6\)。四点置换表示减去平凡项，得到标准表示 \(W\)，其特征标为 \(3,1,-1,0,-1\)。按类大小计算，
\[
	\langle\chi_W,\chi_W\rangle_{S_4} =\frac{1\cdot9+6\cdot1+3\cdot1+8\cdot0+6\cdot1}{24}=1,
\]
所以 \(W\) 不可约。与符号表示张量得到另一不同的次数三表示。

\ProseParagraphBreak
集合 \(\{1,2,3,4\}\) 分成两个无序二元组共有三种方式，记为 \(A=12|34\)、\(B=13|24\)、\(C=14|23\)。换位 \((12)\) 固定 \(A\) 并交换 \(B,C\)；双换位 \((12)(34)\) 固定全部三种配对；三轮换 \((123)\) 将 \(A,C,B\) 依次循环；四轮换 \((1234)\) 交换 \(A,C\) 并固定 \(B\)。像群含有一个换位和一个三轮换，故 \(S_4\) 在这三种配对上的作用给出满同态 \(S_4\to S_3\)，核为
\[
	\bigl\{1,(12)(34),(13)(24),(14)(23)\bigr\}.
\]
拉回复次数二标准表示仍不可约：一个拉回不变子空间也是整个像群 \(S_3\) 的不变子空间。在表列的五种轮换型上，固定配对数依次为 \(3,1,3,0,1\)。配对置换空间中的常向量构成一维平凡直和项，去掉它便将各迹减去一，得到这一行的值 \(2,0,2,-1,0\)。加上平凡、符号两行，平方和 \(1+1+9+9+4=24\)，得到全部不可约，见\cref{b5:tab:s4}。

\begin{table}[htbp]
	\centering
	\caption{\(S_4\) 的特征标表}\label{b5:tab:s4}
	\begin{tabular}{lrrrrr}
		\toprule
		代表&\(1\)&\((12)\)&\((12)(34)\)&\((123)\)&\((1234)\)\\
		类大小&1&6&3&8&6\\
		\midrule
		\(1\)&1&1&1&1&1\\
		\(\mathrm{sgn}\)&1&-1&1&1&-1\\
		\(\chi_W\)&3&1&-1&0&-1\\
		\(\chi_W\mathrm{sgn}\)&3&-1&-1&0&1\\
		\(\psi\)&2&0&2&-1&0\\
		\bottomrule
	\end{tabular}
\end{table}

对 \(A_4\)，三个双换位组成一个类。三轮换在 \(A_4\) 中的中心化子仍为它生成的三阶群，所以每类大小四，八个三轮换分成两个类。一个三轮换与其逆在 \(A_4\) 中不共轭。若 \(a\in A_4\) 满足 \(a(123)a^{-1}=(132)\)，则因 \((23)(132)(23)=(123)\)，有
\[
	(23)a\in C_{S_4}((123))=\langle(123)\rangle.
\]
右边全是偶置换，左边却为奇置换，矛盾。具体地，\((123)\) 的共轭类是
\[
	\{(123),(142),(134),(243)\},
\]
另一类由这四个元素的逆组成。商 \(A_4/V_4\simeq C_3\) 给出三个一维特征标。\(S_4\) 标准表示限制后取值 \(3,-1,0,0\)，范数平方为 \((9+3)/12=1\)，故仍不可约。平方和为 \(3+9=12\)，见\cref{b5:tab:a4}，其中 \(\omega=\exp(2\pi\mathrm i/3)\)。

\begin{table}[htbp]
	\centering
	\caption{\(A_4\) 的特征标表}\label{b5:tab:a4}
	\begin{tabular}{lrrrr}
		\toprule
		代表&\(1\)&\((12)(34)\)&\((123)\)&\((132)\)\\
		类大小&1&3&4&4\\
		\midrule
		\(1\)&1&1&1&1\\
		\(\lambda\)&1&1&\(\omega\)&\(\omega^2\)\\
		\(\lambda^2\)&1&1&\(\omega^2\)&\(\omega\)\\
		\(\theta\)&3&-1&0&0\\
		\bottomrule
	\end{tabular}
\end{table}

\subsection{特征标表的信息与局限}
\label{b5:subsec:0204:03}

从次数列及其平方和可以恢复群阶；列的范数给出中心化子阶，从而给出共轭类大小。一维表示的数目等于交换化群 \(G^{\mathrm{ab}}=G/[G,G]\) 的阶：复群代数 \(\mathbb C[G^{\mathrm{ab}}]\) 的全部不可约表示都一维，再用正则分解即可得到这一计数。特征标表也决定任意两个已知表示的张量积分解，只需逐点相乘后与各不可约行取内积。

\begin{proposition}\label{b5:prop:character-kernel}
	设 \(G\) 为有限群，\(\rho:G\to\operatorname{GL}_{\mathbb C}(V)\) 为 \(d\) 维复表示，其特征标为 \(\chi\)。则
	\[
		\ker\rho=\bigl\{g\in G\mid\chi(g)=d\bigr\}.
	\]
\end{proposition}
\begin{proof}
	一个方向直接取迹。反过来，\(\rho(g)\) 可对角化，特征值 \(\lambda_1,\ldots,\lambda_d\) 均在单位圆上。若其和为正实数 \(d\)，则 \(\sum_i\operatorname{Re}\lambda_i=d\)，而每项至多为一，所以每项实部都为一，即所有 \(\lambda_i=1\)。可对角化性便使 \(\rho(g)=I\)。零表示的核为全群，该公式也成立。
\end{proof}

但普通特征标表不能唯一确定群的同构类型。\Cref{b5:tab:d8-q8}中的两群有相同表，\(D_8\) 却有五个阶二元素，而 \(Q_8\) 只有一个。一个表的列给出共轭类上的函数值，不自动包含各类取平方后落到哪一列的全部信息。若要计算\eqref{b5:eq:symmetric-exterior-character}中的平方类，就须再有这些群结构数据。
